\section{一些集合和公理化的东西}
\subsection{公理性的东西}
\textbf{最小整数公理}$^\clubsuit $:自然数集$\mathbb{N}$的每个非空子集$C$中都含有一个最小的整数.\par
\begin{Theorem}{数学归纳法}
给定一组关于自然数$n\geqslant 1$的命题$S(n)$,假设\\
    (1)\textbf{基础步骤}:$S(1)$成立;\\
    (2)\textbf{归纳步骤}:若$S(n)$成立,则$S(n+1)$也成立.\\
    那么对一切自然数$n\geqslant 1$,命题$S(n)$都成立.
\end{Theorem}
设$X$为一集合,则$X$上的一个关系指的是$X\times X$的一个子集$R$.如果$(x,y)\in R$,记为$xRy$.
\begin{Definition}
如果$R$是$X$上的一个关系,且$\forall x,y,z\in X$满足:\par
(1).\textbf{自反性}:$xRx$;\par
(2).\textbf{对称性}:如果$xRy$,则$yRx$;\par
(3).\textbf{传递性}:如果$xRy$且$yRz$,则$xRz$.\par
那么称$R$为$X$上的一个等价关系.对于等价关系,我们通常记为$\thicksim $.
\end{Definition}
对于任意的$x\in X$,所有与$x$等价的元素组成的集合
$$
[x]=\{y\in X;y\thicksim x\}
$$
称为$x$的等价类.最后考虑$X$所有等价类组成的集合
$$
X/\sim=\{[x];x\in X\}
$$
称为$X$关于等价关系$\thicksim$的商集.显然,商集$X/\sim$是$2^{X}$的子集.\par
在等价关系中有一类特殊的关系,称为序关系.集合$X$上的关系$\prec $是偏序关系,如果它满足:\par
(1).\textbf{自反性}:$x\prec  x$;\par
(2).\textbf{反对称性}:如果$x\prec  y$且$y\prec  x$,则$x=y$;\par
(3).\textbf{传递性}:如果$x\prec  y$且$y\prec  z$,则$x\prec  z$.\par
此时称$(X,\prec)$为一偏序集.如果对于$\forall x,y\in X$,必有$x\prec y$或$y\prec x$,则称$(X,\prec)$为全序集.\par
设$(X,\prec)$为一偏序集,$A$是$X$的一个非空子集,元$s\in X$称为$A$的上界,是指对所有的$a\in A$有$a\prec s$;如果$s\in A$,那么称$s$是$A$的极大元.
\begin{tcolorbox}[colback=KuangNei,colframe=BianKuang,title=\textbf{zorn引理},sharpish corners]
    若非空偏序集$X$的任意全序子集在$X$中都有上界,则$X$必有极大元.
\end{tcolorbox}
有时候我们还会接触到链的概念:偏序集$(X,\prec)$中的链是指子集$Y\subset X$使得$(Y,\prec)$是全序集.因此zorn引理还可以表述为:如果偏序集$X$的任意链在$X$中都有上界,则$X$必有极大元.\par
现在来考虑选择公理,在此之前需要定义选择函数:设$X$是一个集族(元素为非空集合),所谓的选择函数$f$是指定义在$X$上,使得对任意的$A\in X$,$f(A)\in A$.
\begin{tcolorbox}[colback=KuangNei,colframe=BianKuang,title=\textbf{选择公理},sharpish corners]
    设$X$为任一集族,且$\emptyset\notin X$,则存在$f:X\to \bigcup_{A\in X}A$,使得对任意的$A\in X$,都有$f(A)\in A$.
\end{tcolorbox}
\subsection{集合中的一些东西}
\textbf{定义}$^\clubsuit $:如果集合列$\{A_n\}$是单调增加的,那么
$$
\lim_{n\to \infty}A_n=\bigcup_{n=1}^{\infty}A_n
$$
\qquad 如果集合列$\{A_n\}$是单调减少的,那么
$$
\lim_{n\to \infty}A_n=\bigcap_{n=1}^{\infty}A_n
$$
\qquad 我们注意到$\Bigl\{\bigcup_{m=n}^{\infty}A_m\Bigr\}$是关于$n$单调减少的集合列,因此极限存在,于是就得到集合列$\{A_n\}$的上极限
$$
\varlimsup_{n\to\infty}A_n=\bigcap_{n=1}^{\infty}\bigcup_{m=n}^{\infty}A_m
$$
\qquad 类似的$\Bigl\{\bigcap_{m=n}^{\infty}A_m\Bigr\}$是关于$n$单调增加的集合列,因此极限存在,于是就得到集合列$\{A_n\}$的下极限
$$
\varliminf_{n\to\infty}A_n=\bigcup_{n=1}^{\infty}\bigcap_{m=n}^{\infty}A_m
$$
显然对于一般的集合列$\{A_n\}$
$$
\bigcap_{n=1}^{\infty}A_n\subset \varliminf_{n\to\infty}A_n\subset \varlimsup_{n\to\infty}A_n\subset \bigcup_{n=1}^{\infty}A_n
$$
\qquad 设$G$是直线上的开集,如果开区间$(\alpha,\beta)\subset G$,而且端点$\alpha,\beta$不属于$G$,那么称$(\alpha,\beta)$为$G$的一个构成区间.值得注意的是,$G$的任意两个不同的构成区间一定不相交.\par
\textbf{例}$^\spadesuit $:设$A$是实数$(-\infty,\infty)$中一些长度不为零,并且互不相交的开区间构成的集合,即$A=\{(c_\alpha,d_\alpha)|\alpha\in \varLambda ,d_\alpha-c_\alpha>0,\text{并且对于任意的$\alpha,\beta\in \varLambda,\alpha\neq\beta$,$(c_\alpha,d_\alpha)$和$(c_\beta,d_\beta)$不相交}\}$,试证$A$是可列集.\par
\textbf{Proof}:定义映射$f:A\to \mathbb{Q}$,则对每个$(c_\alpha,d_\alpha)$取与$\frac{c_\alpha+d_\alpha}{2}$充分接近的有理数$q_\alpha$,使得有理数$q_\alpha$在$(c_\alpha,d_\alpha)$内.则由$A$中元素是互不相交的开区间可知$f$是单射,故$A$与$\mathbb{Q}$的子集对等,所以$A$一定是有限集或可列集.\par
\textbf{例}$^\spadesuit $:若$U$是实数集的非空开集,则任意$x\in U$,都一定存在构成区间
$(\alpha,\beta)\subseteq U$,使得$x\in(\alpha,\beta).$\par
\textbf{Proof}:对于任意$x\in U$,由$U$是开集可知一定存在某个开区间$(\alpha,\beta)\subseteq U.$ 记$A_x$为满足该性质的开区间全体构成的集合,即$A_x=\{(\alpha,\beta)\mid x\in(\alpha,\beta)$并且$(\alpha,\beta)\subseteq U\}$,明显地,$A_x$不是空集.令
$$
\begin{aligned}
    \alpha_0 &=\inf\{\alpha\mid(\alpha,\beta)\in A_x\}\\
    \beta_0  &=\sup\{\beta\mid(\alpha,\beta)\in A_x\}
\end{aligned}
$$
则$(\alpha_{0},\beta_{0})=\bigcup_{(\alpha,\beta)\subseteq A_{x}}(\alpha,\beta).$\par
明显地,由于$A_x$中的每个开区间都包含$x$,因此$x\in(\alpha_0,\beta_0).$\par
下面先证明$(\alpha_0,\beta_0)\subseteq U$.实际上,对于任意$y\in(\alpha_0,\beta_0)$,不妨设$\alpha_0<$ $y<x.$由于$\alpha_0$是下确界,因此,存在包含$x$的开区间$(\alpha,\beta)$,使得$y>\alpha>\alpha_0$,故$y\in(\alpha,x]\subseteq(\alpha,\beta)\subseteq U$,从而,对于任意 $y\in(\alpha_0,\beta_0)$,都有 $y\in U$,因此$,(\alpha_0,\beta_0)\subseteq U.$\par
再证明$\alpha_0$和$\beta_0$都不属于$U$.用反证法,假设$\alpha_0$属于$U$,则由$U$是开集可知存在开区间$(\alpha^\prime,\beta^{\prime})\subseteq U$,使得$\alpha_0\in(\alpha^\prime,\beta^\prime)$,故$\alpha^\prime<\alpha_0$,并且$(\alpha^\prime,\beta^\prime)\subseteq A_x$,但这与$\alpha_0$是$A_x$中所有开区间的左端点的下确界矛盾.故而,$\alpha_0$不属于$U$.类似可证,$\beta_0$也不属于$U$.\par
综合上可知,($\alpha_0,\beta_0)$是$U$的构成区间.
\begin{tcolorbox}[colback=KuangNei,colframe=BianKuang,title=\textbf{直线上开集的构造},sharpish corners]
    直线上集合$G$是非空开集的充分必要条件是$G$可以表示成有限个或可列个互不相交的构成区间的并集.
\end{tcolorbox}
\textbf{Proof}:充分性显然.只需要证必要性.由上面的例题可以看出$G$是其构成区间的并集,并且其构成区间至多可列.\par
下面的例题一般是经常用到的小结论.\par
\textbf{例}$^\spadesuit $:设$f:[a,b]\to\mathbb{R}$为实函数，则
$$
\bigcup_{n=1}^{\infty}\left\{x\in[a,b]:|f(x)|>\frac{1}{n}\right\}=\{x\in[a,b] |f(x)\neq 0\}
$$
\qquad \textbf{Proof}:证明集合相等首先应当想到证互相包含.即
$$
\begin{aligned}
    &\qquad x_0\in\bigcup_{n=1}^{\infty}\left\{x\in[a,b]:|f(x)|>\frac{1}{n}\right\}\\
    &\Leftrightarrow\exists n_0\in\mathbb{N},s.t.\quad x_0\in\left\{x\in[a,b] :|f(x)|>\frac{1}{n_0}\right\}\\
    &\Leftrightarrow\exists n_0\in\mathbb{N},s.t.\quad |f(x_0)|>\frac{1}{n_0},\text{即}f(x_0)\neq 0\\
    &\Leftrightarrow x_0\in\{x\in[a,b]|f(x)\neq 0\}
\end{aligned}
$$
这就证明了
$$
\bigcup_{n=1}^{\infty}\left\{x\in[a,b] : |f(x)|>\frac{1}{n}\right\}=\{x\in[a,b]|f(x)\neq 0\}
$$
\qquad \textbf{构造不相交的集合列}$^\spadesuit $:设$A_n(a=1,2,...)$为一集合列.令$B_1=A_1,B_n=A_n-\bigcup_{i=1}^{n-1}A_i(n\geqslant 2)$.证明:$B_n(n=1,2,...)$为一个彼此不相交的集列,并且
$$
\bigcup_{i=1}^n A_i=\bigcup_{i=1}^n B_i,\ n=1,2,3\cdots
$$
$$
\bigcup_{i=1}^\infty A_i=\bigcup_{i=1}^\infty B_i
$$
\qquad \textbf{Proof}:显然,$B_1=A_1-\bigcup_{i=1}^{1}A_i$,故$B_1⊂A_1$,故$\bigcup_{i=1}^\infty B_i⊂\bigcup_{i=1}^\infty A_i$.\par
反之,若$x\in\bigcup_{i=1}^\infty A_i,\text{当}\exists x\in A_1$,则$x\in A_1=B_1\subset\bigcup_{i=1}^\infty B_i$;当$x\notin A_i,i=1,2,...,m,x\in A_{m+1}$,则$x\in A_{m+1}-\bigcup_{i=1}^mA_i=B_{m+1}\subset\bigcup_{i=1}^\infty B_i$.因此,$\bigcup_{i=1}^\infty A_i\subset\bigcup_{i=1}^\infty B_i$.\par
综上所述得到
$$
\bigcup_{i=1}^n A_i=\bigcup_{i=1}^n B_i
$$
易见,由$B_i=A_i-\bigcup_{j=1}^{i-1}A_j\subset A_i$,知$\bigcup_{i=1}^\infty B_i\subset\bigcup_{i=1}^\infty A_i$.\par
反之,若$x\in\bigcup_{i=1}^\infty A_i,\text{当}\exists x\in A_1$,则$x\in A_1=B_1\subset\bigcup_{i=1}^\infty B_i;\text{当}x\notin A_i,i=1,2,...,m,x\in A_{m+1}$,则$x\in A_{m+1}-\bigcup_{i=1}^mA_i=B_{m+1}\subset\bigcup_{i=1}^\infty B_i$.因此,$\bigcup_{i=1}^\infty A_i\subset\bigcup_{i=1}^\infty B_i$.\par
综上所述得到
$$
\bigcup_{i=1}^n A_i=\bigcup_{i=1}^n B_i
$$